\section{Introduction}\label{sec:intro}

Forecasters expect artificial intelligence to raise output and, at the same
time, to tilt its composition away from labor and toward capital
\citep{karger2026}. The United States taxes those two kinds of income very
differently, and it runs a transfer system that responds automatically when
wages fall. So the fiscal and household consequences of AI-driven growth
depend not mainly on how much growth arrives, but on which incomes it arrives
in --- and on which parts of the fiscal system a model can see. The Budget
Lab at Yale recently scored AI growth scenarios through a federal individual
income and payroll tax microsimulation and found that a growing capital share
mutes the revenue gains \citep{iselinnunn2026}. Their model, by design, stops
at the tax code: no transfer programs beyond tax credits, no poverty
measurement, no state systems.

This paper runs the same question through the whole system. PolicyEngine-US
is an open-source static microsimulation of United States federal and state
tax and benefit law --- income and payroll taxes, refundable credits, SNAP,
SSI, TANF, WIC and the other major transfer programs, and the Supplemental
Poverty Measure (SPM) --- running on openly published microdata. I use it in
two complementary designs. The first is a mechanism sweep: shift 0--100\% of
positive labor income into capital income, holding modelled market income
fixed, and record what current law does at each step. The second is a
calibrated scenario set that deliberately adopts the Budget Lab's
implementation of the \citet{karger2026} forecaster survey --- Slow, Moderate
and Rapid adoption paths crossed with compressive, proportional and expansive
wage-inequality variants, plus their shares-fixed counterfactuals --- so that
where the two models differ, the difference is model coverage rather than
scenario assumptions. Because they publish their code and their complete
result grid, the comparison can then be made cell by cell on an identical
accounting basis.

Four findings organize the paper. First, the composition tilt, not the growth
rate, drives the fiscal and poverty results. Rapid adoption with factor
shares held at today's levels raises net government revenue by
\genRapidFixedRev B in 2030 and cuts SPM poverty by
\genRapidFixedPovPp pp; the same growth with the forecast tilt keeps
\genKeptHouseholdRapidPct\% of that revenue (\genRapidPropRev B) and turns
the poverty gain into a marginal loss (\genRapidPropPovPp pp). The keep-rate
is basis-dependent, and the paper reports it both ways: on the narrower
federal basis that the Budget Lab measures, the tilt keeps
\genKeptMatchedRapidPct\% here and \genKeptTheirsRapidPct\% in their grid.
The household figure is lower mainly because that basis has no corporate
wedge, which grows with the capital shock; adding state taxes and the
transfer system raises the tilt's dollar cost but barely moves the keep-rate.

Second, the assumption with the least evidence behind it --- whether AI
compresses or widens the wage distribution, a question the forecaster survey
does not elicit and the Budget Lab spans with stylized variants --- decides
the household verdict while leaving the revenue verdict intact. Across the Budget Lab's own
compressive-to-expansive span, Rapid net revenue stays between
\genRapidCompRev B and \genRapidExpRev B, always positive, while the poverty
outcome flips sign: \genRapidCompPovPp pp under compression,
\genRapidExpPovPp pp under expansion, a \genRapidPovSpanPp pp span.

Third, the parts of the fiscal system outside a tax-only model respond
materially, and in both directions. Across those same wage-spread variants,
gross tax collections swing \$\genGrossTaxSwing B, and automatic transfer
responses absorb \genTransferAbsorbPct\% of that swing ---
\$\genBenefitSwing B of it in non-credit programs such as SNAP and SSI that a
tax microsimulation does not carry. State tax codes add
\genRapidPropState B under proportional Rapid growth and double across the
variants, from \genRapidCompState B to \genRapidExpState B. Whether these
channels matter is exactly the kind of question a full-system model exists to
answer.

Fourth, two definitional choices about the tax base move the revenue answer
by more than the choice between Slow, Moderate and Rapid. If AI's incremental
capital income goes unrealized --- retained, accrued, or sheltered inside
retirement accounts --- revenue falls roughly linearly, and below a realized
share of roughly \genRealBreakevenPct\% the household sector loses money on
Rapid AI growth; poverty, meanwhile, moves by
\genRealPovSpanPp~pp across the whole sweep. And whether
taxable retirement distributions count as capital income halves the Rapid
estimate when excluded (\genRapidScopeFull B to \genRapidScopeExcl B).

The reconciliation with the Budget Lab in Section~\ref{sec:reconciliation} is,
to my knowledge, the first cell-level comparison of two independent
microsimulations of AI fiscal scenarios. It is reassuring where the models
share machinery --- payroll responses agree in all nine cells to within
\$\genPayrollMaxCellGap B, two independent wage distributions meeting the
same taxable maximum --- and
informative where they do not: the bridged headline is \genBridgeOursTotal B
against their \genBridgeTheirsTotal B (\genBridgeRatio$\times$), with the gap
concentrated in the income-tax response to the capital shock --- nearly
constant in absolute terms across the wage-spread cells --- consistent with
this paper's \genBaseGapPct\% larger realized-capital base and more
top-concentrated allocation rule.

The paper proceeds as follows. Section~\ref{sec:literature} places the
exercise in the literature. Section~\ref{sec:method} describes the model,
data, and both experiment designs. Section~\ref{sec:results} reports the
sweep and scenario results, including the transfer-system and state channels.
Section~\ref{sec:sensitivity} reports the realization and capital-scope
sensitivities and stability across data builds.
Section~\ref{sec:reconciliation} reconciles with the Budget Lab cell by cell.
Section~\ref{sec:discussion} discusses limitations and next steps.
